
\subsection{Teorema: Linear Speedup}
\subsubsection{Versione tempo}
Se $L \in \TIME(f(n))$ allora $\forall \varepsilon > 0 \quad L \in \TIME(f'(n))$ \\
Dove $f'(n) = \varepsilon f(n) + n + 2$
\begin{proof}
Simuliamo $M=(Q, \Sigma, \delta, s)$ con una macchina $M'=(Q', \Sigma', \delta', s')$. \\
L'idea é simulare $m$ passi di $M$ con $6$ passi di $M'$, scegliendo $m$ opportunamente. \\
Per farlo, occorre utilizzare un numero di stati e di simboli dell'alfabeto esponenziale rispetto a quelli della macchina originale. \\ 
$\Sigma' = \Sigma \cup \Sigma^m$: contiene tutti i simboli dell'alfabeto originale e tutte le $m$-uple di simboli su quell'alfabeto.
Cioé consideriamo $(\sigma_1, \ldots, \sigma_m)$ come un \textit{unico} simbolo. \\

La macchina comincia con una scansione del primo nastro per comprimere l'input: ogni sequenza di $m$ simboli di $x$ viene scritta con un solo simbolo di $\Sigma^m$ sul secondo nastro.\\
Opzionalmente, si puó cancellare $x$ durante la scansione per avere un nastro in piú a disposizione. \\
Notare che se $M$ ha un solo nastro, allora $M'$ ha bisogno di due nastri. In tutti gli altri casi possono avere lo stesso numero di nastri. \\
\begin{itemize}
    \item dallo stato iniziale $s'$ usiamo stati da ${s'}_{\sigma_1}$ a ${s'}_{\sigma_1, \ldots, \sigma_{m-1}}$ per ricordare ogni simbolo incontrato (inclusi i $\sqcup$ se $n$ non é divisibile per $m$).
    \item in uno degli stati che ricordano $m-1$ simboli, leggiamo quello sotto la testina e scriviamo sul secondo nastro $(\sigma_1 \ldots, \sigma_m)$.
    \item Tempo: $m \cdot \ceil{n/m} + 2$  passi sul primo nastro, piú $\ceil{n/m}$ sul secondo.
\end{itemize}

Ora la simulazione comincia usando la stringa del secondo nastro come input. \\
Verranno utilizzati \textit{solo} simboli di $\Sigma^m$ per simulare $m$ passi alla volta. \\

\textbf{Se la testina del nastro $i$ di $M$ si trova sulla cella $l_i$, allora la testina del nastro $i$ di $M'$ si trova sulla cella $\ceil{l_i/m}$.} \\
Di conseguenza, per conoscere la posizione della testina di $M$ sul nastro $i$, basta ricordare solo in quale posizione sta tra le $m$ del simbolo corrente di $M'$. \\
Il numero di stati richiesti per ricordare lo stato di $M$ e la posizione della testina é $|Q| \cdot m^k$. \\

Se $M$ é nello stato $q$ e la testina del suo nastro $i$ si trova sulla cella $l_i$, allora chiamiamo ${l'}_i = (l_i \mod m)$ e diciamo che $M'$ é nello stato $q_{l'_1, \ldots, l'_k}$. .

Simulando $m$ passi di $M$, il cursore di un nastro di $M'$ potrebbe finire in una delle due caselle adiacenti a quella attuale, e non oltre (perché lo spazio é limitato superiormente dal tempo).\\
Per questo, si leggono i simboli a sinistra e a destra del cursore (1 passo a sinistra, 2 a destra e 1 a sinistra tornando alla posizione originaria: 4 in totale). \\
$\delta'$ é definita in modo da applicare direttamente le modifiche al blocco centrale e, eventualmente, uno dei due adiacenti. Questo richiede al massimo 2 passi. \\

In totale, quindi, ogni simulazione di $m$ passi di $M$ richiede $6$ passi di $M'$. \\
$M'$ impiega $6 \cdot \ceil{f(n)/m} + n + 2$ passi per simulare $M$. \\
Imponendo $m = \ceil{6/\varepsilon}$ otteniamo $\varepsilon f(n) + n + 2$.
\end{proof}

Questo teorema ci permette di ignorare costanti moltiplicative nella misurazione della complessitá di tempo.

\subsubsection{Versione spazio}
Se $L \in \SPACE(f(n))$ allora $\forall \varepsilon > 0 \quad L \in \SPACE(f'(n))$ \\
Dove $f'(n) = \varepsilon f(n) + n + 2$

Questo teorema ci permette di ignorare costanti moltiplicative nella misurazione della complessitá di spazio.
